\section{Problematización y Objetivos de Investigación}

\subsection{Problematización Científica}
Los estudios tradicionales de mercado orientados a la inclusión social suelen sufrir de un sesgo crítico: la \textbf{deseabilidad social declarativa}. Cuando se le pregunta directamente a un decisor gerencial si contrataría a una persona con antecedentes o si discrimina a una mujer, la respuesta suele alinearse con la corrección política, ocultando las tasas reales de veto implícito.

Asimismo, existe una desconexión entre la oferta de programas sociales y la estructura de incentivos y riesgos de las empresas según su tamaño (micro, pyme, grande, corporación) y sector económico. Es fundamental precisar estadísticamente qué atributos de la oferta de JxR (acompañamiento, nivel técnico, certificaciones CET/OTEC, reducción de riesgo) mitigan eficazmente la resistencia directiva.

\subsection{Preguntas de Investigación}
\begin{itemize}
    \item \textbf{Pregunta General:} ¿Cuáles son los determinantes causales, barreras organizacionales e incentivos B2B que condicionan el avance de las empresas en la cadena de contratación de personas con antecedentes, y cómo varía esta dinámica según género y segmento empresarial?
    \item \textbf{Preguntas Específicas:}
    \begin{enumerate}
        \item ¿Cuál es la brecha real de retención laboral observada entre hombres y mujeres en los datos administrativos de JxR?
        \item ¿Cuál es la tasa oculta de veto empresarial discriminatorio frente a perfiles con antecedentes penales y su penalización por género, corregida por sesgo de deseabilidad social?
        \item ¿Cómo priorizan los decisores gerenciales las barreras e incentivos de contratación al enfrentar decisiones de elección forzada?
        \item ¿Qué modelo de colaboración (Bolsa Laboral, OTEC o CET) y qué encuadre de mensaje B2B maximiza la intención de reunión y contratación?
    \end{enumerate}
\end{itemize}

\subsection{Objetivos del Estudio}
\begin{itemize}
    \item \textbf{Objetivo General:} Desarrollar un modelo explicativo-predictivo multivariado sobre las barreras e incentivos en la contratación B2B con enfoque de género, proponiendo una estrategia optimizada de prospección y rediseño programático para la Fundación JxR.
    \item \textbf{Objetivos Específicos:}
    \begin{enumerate}
        \item Analizar la línea base de permanencia observada según sexo mediante curvas de Kaplan-Meier sobre la base histórica de JxR.
        \item Mapear cualitativamente la dinámica interna de veto y facilitación mediante tríadas gerenciales y entrevistas en profundidad a mujeres egresadas.
        \item Cuantificar la prevalencia del veto oculto y la brecha de género mediante un experimento de doble lista en una muestra nacional de 300--360 decisores.
        \item Estimar la utilidad jerárquica (HB) de las barreras e incentivos mediante un diseño MaxDiff.
        \item Evaluar causalmente mediante tests monádicos y experimentos B2B los modelos de oferta y mensajes comerciales de la Fundación.
    \end{enumerate}
\end{itemize}